% from ega4-12-fr.tex
\oldpage[IV-3]{174}
\subsection{Propriétés locales des fibres d'un morphisme plat localement de
présentation finie}
\label{subsection:IV.12.1-fr}

\begin{theorem}[12.1.1]
\EGAReaderTarget{ega4.c04.thm.12-1-1}{IV.12.1.1-fr}
Soient $f:X\to Y$ un morphisme localement de présentation finie, $\mathcal F$
un $\mathcal O_X$-Module $f$-plat et de présentation finie, $\Phi$ une partie
finie de $\mathbf Z\cup\{\pm\infty\}$, $k$ un entier. Les parties suivantes de
$X$ sont ouvertes~:
\begin{enumerate}[label=(\roman*)]
  \item L'ensemble des $x\in X$ tels que les dimensions des cycles premiers
  associés à $\mathcal F_{f(x)}$ et contenant $x$ soient des éléments de
  $\Phi$.
  \item L'ensemble des $x\in X$ tels que les cycles premiers associés à
  $\mathcal F_{f(x)}$ contenant $x$ aient tous la dimension $k$.
  \item L'ensemble des $x\in X$ tels que $x$ n'appartienne à aucun cycle
  premier immergé associé à $\mathcal F_{f(x)}$.
  \item L'ensemble des $x\in X$ tels que $\mathcal F_{f(x)}$ soit
  équidimensionnel au point $x$ et possède la propriété $(S_k)$ au point $x$
  (5.7.2).
  \item L'ensemble des $x\in X$ tels que
  $\operatorname{coprof}((\mathcal F_{f(x)})_x)\leq k$
  ($\mathbf 0$, 16.4.9).
  \item L'ensemble des $x\in X$ tels que $\mathcal F_{f(x)}$ soit un
  $\mathcal O_{X_{f(x)}}$-Module de Cohen-Macaulay au point $x$ (5.7.1).
  \item L'ensemble des $x\in X$ tels que $\mathcal F_{f(x)}$ soit
  géométriquement réduit au point $x$ (4.6.22).
  \item L'ensemble des $x\in X$ tels que $\mathcal F_{f(x)}$ soit
  géométriquement ponctuellement intègre au point $x$ (4.6.22).
\end{enumerate}
\end{theorem}

\begin{proof}
Les questions étant locales sur $X$, on se ramène d'abord au cas où
$X=\operatorname{Spec}(B)$ et $Y=\operatorname{Spec}(A)$ sont affines, $f$ un
morphisme de présentation finie. Il y a alors un sous-anneau noethérien $A_0$
de $A$, un $A_0$-préschéma de type fini $X_0$, un
$\mathcal O_{X_0}$-Module cohérent $\mathcal F_0$ tel que
$\mathcal F_0\otimes_{A_0}A$ soit isomorphe à $\mathcal F$ (8.9.1)~; en
outre, en vertu de (11.2.6), on peut supposer que $\mathcal F_0$ est
$Y_0$-plat (avec $Y_0=\operatorname{Spec}(A_0)$). Si $h:X\to X_0$ est la
projection canonique, l'ensemble $E$ des points $x\in X$ où l'une des
propriétés (i) à (viii) est vérifiée est égal à $h^{-1}(E_0)$, où $E_0$ est
l'ensemble des $x_0\in X_0$ où la propriété correspondante pour $\mathcal F_0$
est vérifiée~: cela résulte, pour les propriétés (i) à (iii), de (4.2.7)~; pour
les propriétés (iv) à (vi), de (6.7.1)~; pour les propriétés (vii) et (viii),
de (4.7.11).

On est ainsi ramené au cas où $A$ et $B$ sont noethériens, $f$ de type fini,
et $\mathcal F=\widetilde M$, où $M$ est un $B$-module de type fini. En vertu
de (9.9.2), on sait que l'ensemble $E$ est constructible pour toutes les
propriétés envisagées, et il reste dans chaque cas l'étape C) de (12.0.2), où
l'on doit prouver que $E$ est stable par générisation.

\begin{env}[12.1.1.1]
\label{IV.12.1.1.1-fr}
Commençons par le cas (iii) qui est le plus simple. Soit $x'\ne x$ une
générisation de $x$ dans $X$. Posons $y=f(x)$, $y'=f(x')$~; il existe un
spectre d'anneau de valuation discrète $Y_1$ et un morphisme $u:Y_1\to X$ tels
que, si $y_1$ et $y'_1$ sont le point fermé et le point générique de $Y_1$, on
ait $u(y_1)=x$ et $u(y'_1)=x'$ ($\mathbf{II}$, 7.1.7)~; posons
$g=f\circ u:Y_1\to Y$~; si $X_1=X\times_YY_1$, il existe une $Y_1$-section
$u_1:Y_1\to X_1$ telle que $u=g_1\circ u_1$, où $g_1:X_1\to X$ est la
projection canonique. Si $x_1=u_1(y_1)$, $x'_1=u_1(y'_1)$, on a donc
$g_1(x_1)=x$ et $g_1(x'_1)=x'$, et $x'_1$ est une générisation de $x_1$.
Appliquant de nouveau (4.2.7), on voit
